\documentclass[aspectratio=169]{beamer}
\usetheme{Madrid}
\usecolortheme{seahorse}

\usepackage{graphicx}
\usepackage{booktabs}
\usepackage{amsmath}
\usepackage{amssymb}
\usepackage{tabularx}
\usepackage{array}
\usepackage{ragged2e}

\title[Chapter 6: Samples and Observational Studies]{\textbf{Chapter 6: Samples and Observational Studies}}
\subtitle{The Practice of Statistics in the Life Sciences \\ by Brigitte Baldi and David S. Moore}
\author{Instructor: Ali Raisolsadat}
\date{Part III: Data Generation}

\setbeamertemplate{navigation symbols}{}
\setbeamertemplate{itemize items}[circle]
\setbeamertemplate{enumerate items}[default]
\graphicspath{{./}{figures/}}

\newcommand{\lectureimage}[2][0.90\linewidth]{%
  \IfFileExists{#2}{%
    \includegraphics[width=#1]{#2}%
  }{%
    \fbox{\parbox[c][0.36\textheight][c]{0.88\linewidth}{%
      \centering\small
      Figure file: \texttt{\detokenize{#2}}\\[0.15cm]
      Generate this figure by knitting the revised Chapter 6 RMarkdown notes.
    }}%
  }%
}

\begin{document}

\begin{frame}
  \titlepage
\end{frame}

\begin{frame}{Why This Chapter Matters}
Statistical methods are useful only when the data have been collected in a way that matches the scientific question.

\vspace{0.15cm}
Before a mean, proportion, confidence interval or test is interpreted, the study design should make clear
\begin{itemize}
  \item who or what is being studied
  \item what variables are being measured
  \item how the observations are obtained
\end{itemize}

The quality of a statistical conclusion depends on the quality of the data-generation process.
\end{frame}
%
%\begin{frame}{Chapter Roadmap}
%This chapter develops the main ideas in a sequence that matches the way studies are planned.
%\begin{enumerate}
%  \item Observational studies and experiments
%  \item Lurking variables and confounding
%  \item Populations, samples and parameters
%  \item Sampling designs and probability sampling
%  \item Bias in sample surveys
%  \item Case-control and cohort studies
%  \item Matching conclusions to study design
%\end{enumerate}
%\end{frame}

\section{Observational Studies and Experiments}

\begin{frame}{Gathering Data}
Many scientific questions concern populations that are too large to study completely.

\vspace{0.15cm}
A health researcher may want to estimate a proportion for adults in a country. A wildlife biologist may want to learn about fish in a river system. A laboratory may want to understand a biological process from a manageable set of specimens.

\vspace{0.15cm}
In each setting, the main issue is not simply how much data are collected. The main issue is whether the data were collected through a design that supports the intended conclusion.
\end{frame}

\begin{frame}{Study Design Comes Before Statistical Analysis}
A numerical method cannot repair a study design that systematically excludes important groups, measures the wrong variable or mixes the effect of an explanatory variable with other factors.

\vspace{0.15cm}
Statistical reasoning therefore begins with a clear description of how the data were generated.
\end{frame}

\begin{frame}{Observation versus Experiment}
Two broad approaches are used to gather data in the life sciences.

\begin{block}{Observational Study}
An observational study records information without deliberately assigning the main explanatory condition being studied.
\end{block}

\begin{block}{Experiment}
An experiment deliberately imposes one or more treatments on experimental units and compares their responses.
\end{block}
\end{frame}

\begin{frame}{Why the Distinction Matters}
The distinction matters because observational associations can have several explanations.

\vspace{0.15cm}
Individuals who differ in the explanatory variable may also differ in age, behaviour, environment, socioeconomic conditions or other characteristics.

\vspace{0.15cm}
A well-designed randomized experiment helps separate the treatment effect from those other variables, which strengthens causal interpretation.
\end{frame}

\begin{frame}{When Are Observational Studies Needed}
Observational studies remain essential when an experiment would be unethical, impossible or scientifically inappropriate.

\vspace{0.15cm}
For example, researchers should not randomly assign people to smoke cigarettes in order to study long-term health effects.

\vspace{0.15cm}
In these situations, strong scientific conclusions still depend on careful design, subject-matter knowledge and attention to alternative explanations.
\end{frame}

\begin{frame}{Random Sampling and Random Assignment}
These two ideas solve different problems.

\begin{itemize}
  \item \textbf{Random sampling} helps support generalization from a sample to a population
  \item \textbf{Random assignment} helps support causal comparisons between treatments
\end{itemize}

A study can use one, both or neither of these forms of randomness.
\end{frame}

\section{Confounding and Lurking Variables}

\begin{frame}{Lurking Variables and Confounding}
The main difficulty in an observational study is that the explanatory variable can be related to other variables that also affect the response.

\begin{block}{Lurking Variable}
A lurking variable is a variable that is not included as an explanatory or response variable in the analysis but may help explain the observed relationship.
\end{block}

\begin{block}{Confounding}
Confounding occurs when the effects of two or more explanatory variables on a response cannot be separated clearly from the available data.
\end{block}
\end{frame}

\begin{frame}{How the Ideas Differ}
A lurking variable is a possible source of an alternative explanation.

\vspace{0.15cm}
Confounding describes the resulting difficulty in separating effects.

\vspace{0.15cm}
The ideas are related, but they are not identical.
\end{frame}

\begin{frame}{Reducing Confounding}
Confounding is a major concern in observational research because naturally occurring groups can differ in many ways.

\vspace{0.15cm}
Researchers may try to reduce confounding through methods such as restriction, matching, stratification or statistical adjustment.

\vspace{0.15cm}
These methods can reduce the problem, but unmeasured differences may still remain.
\end{frame}

\begin{frame}{Example 6.2: Wine, Beer, or Spirits?}
\begin{block}{Question}
Some observational studies have reported that health outcomes differ among moderate drinkers who prefer wine, beer or spirits.

Does an observed difference in health establish that beverage type itself causes the difference?
\end{block}
\end{frame}

\begin{frame}{Example 6.2: Explanation}
Beverage preference can be related to many other characteristics.

\vspace{0.15cm}
Groups that prefer different beverages may also differ in diet, smoking behaviour, exercise, income, education and access to health care.

\vspace{0.15cm}
If these characteristics also relate to the health outcome, the effect of beverage preference cannot be separated cleanly from the effects of the other variables.
\end{frame}

\begin{frame}{Example 6.2: Confounding Diagram}
  \centering
  \lectureimage[0.8\linewidth]{fig_ch6_confounding.png}
\end{frame}

\begin{frame}{Example 6.2: Solution}
The observational association does not by itself establish a causal effect of beverage type.

\vspace{0.15cm}
Potential confounders provide alternative explanations for the observed difference.

\vspace{0.15cm}
A scientifically appropriate conclusion is that beverage preference and health are associated in the observed data after considering the study design. A stronger causal claim would require stronger evidence and careful control of competing explanations.
\end{frame}

\section{Sampling}

\begin{frame}{Sampling}
A sample is a subset of a larger population.

\vspace{0.15cm}
Sampling is necessary because a complete census is often too expensive, too slow or impossible.

\vspace{0.15cm}
The goal is not simply to obtain many observations. The goal is to obtain observations through a design that supports the intended generalization.
\end{frame}

\begin{frame}{Core Terminology}
Several terms should be distinguished before a sample is selected.

\begin{itemize}
  \item \textbf{Population:} the full group of individuals or units about which information is desired
  \item \textbf{Sample:} the subset of individuals or units from which data are actually collected
  \item \textbf{Sampling frame:} the practical list or mechanism used to identify units that can be selected
  \item \textbf{Sampling design:} the planned method used to choose the sample
\end{itemize}
\end{frame}

\begin{frame}{Parameter and Statistic}
Two additional terms connect sampling to later inference.

\begin{block}{Parameter}
A numerical characteristic of a population
\end{block}

\begin{block}{Statistic}
A numerical characteristic calculated from a sample
\end{block}

Statistical inference uses sample statistics to learn about unknown population parameters.
\end{frame}

\begin{frame}{Why a Representative Sample Matters}
A spoonful from a well-mixed pot of soup can give useful information about the flavour of the pot because the ingredients are distributed throughout the mixture.

\vspace{0.15cm}
The analogy should not be interpreted to mean that one spoonful represents the pot perfectly. Small samples can still vary by chance.

\vspace{0.15cm}
Human populations and ecological systems are often much more heterogeneous than soup, so a deliberate sampling design is needed.
\end{frame}

\begin{frame}{Planning a Sample}
A useful sampling plan begins with several questions.
\begin{enumerate}
  \item What is the target population
  \item What variables will be measured and how will they be defined
  \item What sampling frame will be used
  \item What sampling method will select the units
  \item How will missing data and nonresponse be handled
\end{enumerate}
\end{frame}

\begin{frame}{Defining a Population and a Variable}
\begin{block}{Question}
A university health centre wants to estimate the proportion of first-year students who usually sleep at least 7 hours per night.

What is the population and what is the main variable?
\end{block}
\end{frame}

\begin{frame}{Defining a Population and a Variable: Solution}
The population is all first-year students at that university.

\vspace{0.15cm}
The main variable is whether a student usually sleeps at least 7 hours per night.

\vspace{0.15cm}
This variable can be recorded as a yes-or-no categorical variable, or the number of hours slept can be recorded as a quantitative variable and then classified if needed.
\end{frame}

\begin{frame}{Generalizing from a Sample}
A scientifically sensible generalization requires more than a large sample.

\vspace{0.15cm}
The sample must be connected to the target population through a design that gives relevant individuals a fair chance of selection.

\vspace{0.15cm}
If important groups are systematically missed, the result can be biased even when the sample size is large.
\end{frame}

\section{Sampling Designs}

\begin{frame}{Sampling Variability and Bias}
Two ideas are central in sampling.

\begin{block}{Sampling Variability}
The natural sample-to-sample variation that occurs because different random samples usually produce different statistics
\end{block}

\begin{block}{Bias}
A systematic tendency for a sampling method or measurement process to miss the true population value in a particular direction
\end{block}
\end{frame}

\begin{frame}{Archery Analogy for Sampling Designs}
  \centering
  \lectureimage[0.5\linewidth]{fig_ch6_archery.png}
\end{frame}

\begin{frame}{Reading the Archery Analogy}
In the archery figure, each point represents an estimate from a repeated sample.

\vspace{0.15cm}
The centre of the target represents the true population value.

\vspace{0.15cm}
The figure illustrates four situations.
\begin{itemize}
  \item Small sample and high variability
  \item Large sample and high variability
  \item Small sample and low variability
  \item Biased sampling that is systematically off-centre
\end{itemize}
\end{frame}

\begin{frame}{Key Lesson from the Archery Figure}
Larger samples often reduce sampling variability.

\vspace{0.15cm}
Larger samples do not automatically remove systematic bias.

\vspace{0.15cm}
A method that consistently misses important parts of the population can remain wrong even when the sample becomes very large.
\end{frame}

\begin{frame}{Poor Sampling Designs}
Not all samples are created in ways that support reliable inference.

\begin{itemize}
  \item Convenience samples are based on easy availability
  \item Voluntary response samples are formed by individuals who choose themselves to respond
\end{itemize}

Both approaches can produce strong bias because participation is not controlled by a suitable probability mechanism.
\end{frame}

\begin{frame}{Convenience Samples}
A convenience sample consists of units that are easiest to reach.

\vspace{0.15cm}
Examples include surveying people in one classroom, taking the first specimens that arrive in a lab or interviewing shoppers who happen to pass a kiosk.

\vspace{0.15cm}
These samples may be useful for preliminary work, but they usually do not support confident generalization to a broader population.
\end{frame}

\begin{frame}{Voluntary Response Samples}
A voluntary response sample is formed when people choose for themselves whether to participate.

\vspace{0.15cm}
Examples include internet polls and call-in surveys.

\vspace{0.15cm}
People with especially strong opinions are often more likely to respond, which can create severe bias.
\end{frame}

\begin{frame}{Probability Samples}
A probability sample uses a known chance mechanism to select units.

\vspace{0.15cm}
These designs create a clearer connection between the sample and the population.

\vspace{0.15cm}
The simplest example is a simple random sample.
\end{frame}

\begin{frame}{Simple Random Samples}
In a simple random sample of size $n$, every possible sample of size $n$ has the same chance of being selected.

\vspace{0.15cm}
This design helps protect against selection bias and provides a foundation for statistical inference.
\end{frame}

\begin{frame}{Other Probability Designs}
Other common probability designs include
\begin{itemize}
  \item \textbf{Stratified samples}, which sample separately within identified groups
  \item \textbf{Cluster samples}, which randomly select groups or clusters, then observe units within them
  \item \textbf{Multistage samples}, which apply random selection in several stages
\end{itemize}

These designs are often more practical than a single simple random sample in large populations.
\end{frame}

\begin{frame}{Random Sampling versus Random Assignment}
Random sampling and random assignment both use chance, but they serve different scientific purposes.

\begin{itemize}
  \item Random sampling supports generalization to a population
  \item Random assignment supports causal comparison of treatments
\end{itemize}

Confusing these ideas is a common first-year mistake, so it is important to keep their purposes separate.
\end{frame}

\section{Sample Surveys}

\begin{frame}{Sample Surveys}
A sample survey is a study that gathers information from a sample in order to describe a population.

\vspace{0.15cm}
Even a large survey can be misleading if some groups are missed, many selected individuals do not respond or the questions are poorly designed.
\end{frame}

\begin{frame}{Undercoverage}
Undercoverage occurs when some groups in the population are left out of the sampling frame or are given little chance of selection.

\vspace{0.15cm}
A survey that samples only landline telephone numbers, for example, can underrepresent people who rely mainly on cell phones.
\end{frame}

\begin{frame}{Example 6.8: The Cell Phone Problem}
\begin{block}{Question}
A health survey uses only landline telephone numbers to sample adults. Why can this design produce undercoverage bias?
\end{block}
\end{frame}

\begin{frame}{Example 6.8: Solution}
Adults who rely mainly on cell phones may be missed or underrepresented.

\vspace{0.15cm}
If cell-phone-only adults differ from landline users on the characteristic being measured, the sample statistic can be systematically different from the population value.

\vspace{0.15cm}
The problem is not simply that the sample is smaller. The problem is that the sampling frame excludes part of the target population.
\end{frame}

\begin{frame}{Nonresponse Bias}
Nonresponse bias occurs when selected individuals do not participate and the nonrespondents differ in important ways from those who do respond.

\vspace{0.15cm}
Weighting or follow-up efforts may reduce the problem, but they do not automatically remove it.
\end{frame}

\begin{frame}{Response Bias}
Response bias occurs when the answers given are systematically inaccurate.

\vspace{0.15cm}
Sensitive topics, interviewer effects, question wording and social pressure can all affect how people answer.
\end{frame}

\begin{frame}{Wording and Order of Questions}
The wording and order of survey questions can influence the responses.

\vspace{0.15cm}
Leading wording can push respondents toward a particular answer. Earlier questions can also shape how later questions are interpreted.
\end{frame}

\begin{frame}{Example 6.10: Question Order and Reported Happiness}
\begin{block}{Question}
A survey first asks people about recent dating experiences and then asks how happy they are overall. Why might the answer to the second question be affected by the first?
\end{block}
\end{frame}

\begin{frame}{Example 6.10: Solution}
The first question can bring a specific aspect of life to mind and make it unusually prominent when the respondent answers the second question.

\vspace{0.15cm}
The overall happiness response may therefore reflect the temporary influence of the earlier question rather than the person\textquotesingle s broader long-run judgment.
\end{frame}

\begin{frame}{A Survey Quality Checklist}
When evaluating a survey, useful questions include
\begin{itemize}
  \item Who is the target population
  \item What is the sampling frame
  \item How were units selected
  \item What was the response rate
  \item Could wording or ordering influence the answers
\end{itemize}
\end{frame}

\section{Cohort and Case-Control Studies}

\begin{frame}{Cohort and Case-Control Studies}
Case-control and cohort studies are observational designs used to study relationships between exposures and outcomes.

\vspace{0.15cm}
They are especially important when an experiment would be unethical, impractical or too slow.
\end{frame}

\begin{frame}{Case-Control Studies}
A case-control study selects individuals according to their outcome status.

\vspace{0.15cm}
The cases have the outcome of interest. The controls do not and should ideally come from the same source population.

\vspace{0.15cm}
This design is especially useful for rare outcomes because researchers do not need to wait for enough cases to appear in a very large cohort.
\end{frame}

\begin{frame}{Case-Control Studies: What They Can and Cannot Do}
Case-control studies can efficiently examine whether previous exposures were more common among cases than among controls.

\vspace{0.15cm}
They cannot directly estimate population prevalence or incidence from the researcher-chosen case-control ratio.
\end{frame}

\begin{frame}{Real-World Case Study: Reye's Syndrome and Aspirin Exposure}
\begin{block}{Question}
Why is a case-control design efficient for a rare outcome such as Reye\textquotesingle s syndrome? What population quantity cannot be estimated directly from the researcher-chosen case-control ratio?
\end{block}
\end{frame}

\begin{frame}{Real-World Case Study: Explanation}
The design begins by locating cases, so it does not require a huge random sample while waiting for a rare outcome to occur.

\vspace{0.15cm}
Controls provide a comparison group for examining whether previous exposures were more common among cases.
\end{frame}

\begin{frame}{Real-World Case Study: Solution}
The design is efficient because it concentrates data collection on individuals who are informative about a rare outcome.

\vspace{0.15cm}
It cannot directly estimate population prevalence or incidence from the case-control sampling ratio.

\vspace{0.15cm}
It can, however, estimate measures of association between exposure and outcome when the cases and controls are selected appropriately.
\end{frame}

\begin{frame}{Cohort Studies}
A cohort study identifies a group of individuals according to exposures or characteristics, then compares the occurrence of outcomes across those groups.

\vspace{0.15cm}
A prospective cohort collects exposure information before the outcomes occur and follows participants forward through time.
\end{frame}

\begin{frame}{Strengths and Limitations of Cohort Studies}
Cohort studies can estimate incidence within the observed cohort, compare outcome rates across exposure groups and study several outcomes from the same exposures.

\vspace{0.15cm}
Their main practical disadvantages include time, cost and loss to follow-up.

\vspace{0.15cm}
Their main statistical limitation is that exposure groups are not created by random assignment, so confounding can remain.
\end{frame}

\begin{frame}{Real-World Case Study: The British Doctors Study}
\begin{block}{Question}
What information can a cohort design provide in a study of smoking and later health outcomes? What limitations remain?
\end{block}
\end{frame}

\begin{frame}{Real-World Case Study: Explanation}
A prospective cohort records exposure information before many outcomes occur, which clarifies the time order between exposure and outcome.

\vspace{0.15cm}
Because the cohort is followed over time, researchers can estimate rates of new outcomes and compare those rates across exposure groups.
\end{frame}

\begin{frame}{Real-World Case Study: Solution}
A cohort study can estimate incidence within the observed cohort, compare outcome rates across exposure groups and study several outcomes from the same exposures.

\vspace{0.15cm}
Confounding can still remain because the exposure groups were not formed by random assignment.

\vspace{0.15cm}
Attrition can also create selection bias if loss to follow-up is related to both exposure and outcome.
\end{frame}

\begin{frame}{Comparing the Observational Designs}
\small
\begin{table}
\centering
\begin{tabularx}{\textwidth}{>{\RaggedRight\arraybackslash}p{0.18\textwidth} >{\RaggedRight\arraybackslash}p{0.24\textwidth} >{\RaggedRight\arraybackslash}p{0.26\textwidth} X}
\toprule
Feature & Sample survey & Case-control study & Cohort study \\
\midrule
Typical starting point & A sample from a population & Outcome status & A defined cohort and exposure information \\
Main direction & Usually a population snapshot & Often looks backward for exposure & Often follows outcomes forward \\
Useful for rare outcomes & Usually inefficient & Especially useful & Can require a very large cohort \\
Can estimate prevalence directly & Yes, with appropriate sampling and measurement & Not from the chosen case-control ratio & Sometimes, depending on the cohort and timing \\
Can estimate incidence directly & Usually not from one cross-sectional survey & No & Yes, within a properly followed cohort \\
Main concerns & Coverage, nonresponse and response bias & Selection, recall and confounding & Confounding, attrition and cost \\
\bottomrule
\end{tabularx}
\end{table}
\end{frame}

\begin{frame}{From Observational Studies to Experiments}
Observational studies can provide strong and scientifically important evidence, especially when experiments are impossible or unethical.

\vspace{0.15cm}
Their interpretation must still consider alternative explanations such as confounding and selection bias.

\vspace{0.15cm}
A randomized experiment addresses a different problem by assigning treatments through chance, which strengthens causal interpretation.
\end{frame}

\section{Summary}

\begin{frame}{Chapter Summary: Core Ideas}
Data collection determines what a statistical analysis can support.

\vspace{0.12cm}
An observational study measures naturally occurring variables without deliberately assigning the main explanatory condition. An experiment deliberately assigns treatments.

\vspace{0.12cm}
Probability sampling connects a sample to a population. Increasing sample size can reduce sampling variability, but it does not automatically remove bias.

\vspace{0.12cm}
Case-control and cohort studies are important observational designs with different strengths and limitations.
\end{frame}

\begin{frame}{A Consistent Study-Design Workflow}
A first-year analysis of study design becomes easier when the same questions are asked in a consistent order.
\begin{enumerate}
  \item What is the scientific question
  \item What is the target population
  \item What are the explanatory and response variables
  \item Is the design observational or experimental
  \item If sampling is used, what frame and selection method were used
  \item Could bias or confounding affect the conclusion
  \item Does the final conclusion match what the design can support
\end{enumerate}
\end{frame}

\end{document}
